% ============================================================
%  Préambule du mémoire — modèle inspiré des mémoires de Licence IFRI/UAC
%  Moteur : pdfLaTeX. Bibliographie : manuelle pour l'instant (voir biblio/).
% ============================================================

\usepackage[utf8]{inputenc}
\usepackage[T1]{fontenc}
\usepackage[french]{babel}
\usepackage{mathpazo}                 % Palatino, comme les mémoires IFRI
\usepackage[scaled=0.92]{helvet}      % Sans-serif pour la page de garde
\usepackage{microtype}
\usepackage[a4paper,left=2.5cm,right=2.5cm,top=2.7cm,bottom=2.7cm,headheight=15pt]{geometry}
\usepackage{setspace}
\setstretch{1.12}
\usepackage{xcolor}
\usepackage{graphicx}
\usepackage{booktabs}
\usepackage{array}
\usepackage{tabularx}
\usepackage{enumitem}
\usepackage{amsmath}
\usepackage{tikz}
\usepackage{caption}
\usepackage{float}
\usepackage{fancyhdr}
\usepackage{titlesec}
\usepackage{etoc}
\usepackage{url}

\definecolor{lienbleu}{RGB}{66,103,200}
\usepackage[colorlinks=true,linkcolor=lienbleu,citecolor=lienbleu,urlcolor=lienbleu,
            pdfauthor={\auteur},pdftitle={\titrememoire}]{hyperref}

% --- Légendes : "FIGURE 2.1 : ..." (sous la figure), "TABLEAU 2.1 : ..." (au-dessus du tableau)
\addto\captionsfrench{\renewcommand{\tablename}{Tableau}\renewcommand{\figurename}{Figure}}
\captionsetup{labelfont=sc,labelsep=colon,font=small}
\captionsetup[table]{position=above}
\captionsetup[figure]{position=below}
\addto\captionsfrench{\renewcommand{\listfigurename}{Liste des figures}\renewcommand{\listtablename}{Liste des tableaux}}

% --- Puces : tirets pour les listes de mémoire, listes compactes
\setlist[itemize]{label=\textbullet,itemsep=0.4em,topsep=0.4em}
\setlist[enumerate]{itemsep=0.4em,topsep=0.4em}

% --- Marqueurs de travail (visibles en rouge dans le PDF tant qu'ils ne sont pas traités)
\newcommand{\asourcer}{\textcolor{red}{\textbf{[SOURCE À FOURNIR]}}}
\newcommand{\adonner}[1]{\textcolor{red}{\textbf{[#1]}}}
\newcommand{\aredigerr}{\par\noindent\textcolor{gray}{\emph{À rédiger.}}\par}

% --- Titres
\titleformat{\section}{\normalfont\Large\bfseries}{\thesection}{0.8em}{}
\titleformat{\subsection}{\normalfont\large\bfseries}{\thesubsection}{0.8em}{}
\titleformat{\subsubsection}{\normalfont\normalsize\bfseries}{\thesubsubsection}{0.8em}{}
\titlespacing*{\section}{0pt}{2.2ex plus 1ex}{1.2ex}
\titlespacing*{\subsection}{0pt}{1.8ex plus 1ex}{1ex}
\titlespacing*{\subsubsection}{0pt}{1.4ex plus 1ex}{0.8ex}

% Cadre « Chapitre N » (inspiré des mémoires IFRI)
\newcommand{\cadrechapitre}{%
  \noindent\begin{tikzpicture}
    \def\h{1.55}
    \draw[line width=0.4pt] (0,0.05) -- (0,\h) -- (3.0,\h);
    \node[anchor=south west,font=\sffamily\large,inner sep=0pt] at (0.05,0.12) {Chapitre};
    \node[anchor=south west,font=\fontsize{52}{52}\selectfont,inner sep=0pt] at (2.35,0.00) {\thechapter};
    \draw[line width=0.4pt] (3.55,\h) -- (\linewidth,\h) -- (\linewidth,0.05) -- (3.55,0.05);
  \end{tikzpicture}}

\titleformat{\chapter}[display]{\normalfont}{\cadrechapitre}{1.2em}{\fontsize{26}{32}\selectfont}
\titleformat{name=\chapter,numberless}[display]{\normalfont}{}{0pt}{\fontsize{26}{32}\selectfont}
\titlespacing*{\chapter}{0pt}{20pt}{28pt}

% Chapitre non numéroté, ajouté au sommaire et à la table des matières
\newcommand{\chapnn}[1]{%
  \chapter*{#1}%
  \addcontentsline{toc}{chapter}{#1}%
  \markboth{#1}{#1}}

% --- En-têtes et pieds de page
\pagestyle{fancy}
\fancyhf{}
\fancyhead[L]{\scriptsize\nouppercase{\leftmark}}
\fancyhead[R]{\scriptsize\nouppercase{\rightmark}}
\fancyfoot[C]{\thepage}
\renewcommand{\headrulewidth}{0.4pt}
\renewcommand{\chaptermark}[1]{\markboth{Chapitre \thechapter.\ #1}{}}
\renewcommand{\sectionmark}[1]{\markright{\thesection.\ #1}}
\fancypagestyle{plain}{\fancyhf{}\fancyfoot[C]{\thepage}\renewcommand{\headrulewidth}{0pt}}

% --- Sommaire (chapitres seulement) et table des matières complète
\newcommand{\sommaire}{%
  \etocsettocstyle{\chapter*{Sommaire}\markboth{}{}}{}%
  \etocsetnexttocdepth{chapter}%
  \tableofcontents}
\newcommand{\tabledesmatieres}{%
  \etocsettocstyle{\chapter*{Table des matières}\markboth{}{}}{}%
  \etocsetnexttocdepth{subsubsection}%
  \tableofcontents}

% --- Entrées de listes (acronymes, glossaire)
\newcommand{\acro}[2]{\par\noindent\textbf{#1 :}\hspace{0.5em}#2\par\vspace{0.6em}}
\newcommand{\gloss}[2]{\par\noindent\textbf{#1 :}\hspace{0.5em}#2\par\vspace{0.8em}}
